\section{Discussion}
\label{sec:discussion}

\paragraph{Three behavioral modes under sustained pressure.} The per-turn slope analysis on MMLU-Pro and GPQA-Diamond (\autoref{tab:rq3}) together with the consensus-loss analysis on HLE identify three qualitatively distinct modes that describe how models respond to escalating adversarial pressure. \emph{Resisters} (Claude Haiku, Claude Sonnet, GPT-5.4) show negative slopes across $T_1$--$T_6$: their initial brittleness at the first pressure turn does not compound,
and many flips are reversed by later self-corrections. \emph{Capitulators} (GPT-5.4 Mini, GPT-5.4 Nano) show the opposite pattern: each additional turn raises flip probability further, and once the model yields it stays yielded. \emph{Fragmenters}, observed distinctly on HLE for Gemini 3.5 Flash, follow neither pattern: the model neither defends nor commits, but instead loses coherent self-consistency, producing per-turn answer distributions with no single majority answer on $72.2\%$ of $T_0$-correct queries. The three modes are not artifacts of model scale: both resisters and capitulators include models from the same family at adjacent capability tiers.

\paragraph{Why the distinction matters for deployment.} Each mode implies a different threat surface. \emph{Resisters} are safest under extended adversarial dialogue: their initial vulnerability stabilizes or recovers, so the operational risk is concentrated at $T_1$ and decays. \emph{Capitulators} become \emph{more} dangerous as
conversations lengthen, inverting the intuition that more deliberation produces safer outputs; for these models, each additional pressure
turn is an additional opportunity to lock in a wrong answer. \emph{Fragmenters} are the most operationally unreliable: even when they do not capitulate to a specific wrong answer, they cease producing a coherent reply, which is a distinct failure mode relative to the user who needs a usable response. A single-turn benchmark detects none of these distinctions, since all three modes can yield the same
$T_1$ flip rate while diverging sharply over $T_2$--$T_6$. Evaluations designed for deployment safety should report the per-turn slope and
the consensus-loss rate alongside any aggregate flip rate.

\paragraph{Pressure as an amplifier, not an independent driver.} The significant positive $\hat\beta_{PH}$ interaction in \autoref{fig:rq2_coefs} formalizes a property visible throughout our analyses: pressure does not produce uniform increases in flip probability regardless of prior belief state, but instead amplifies pre-existing uncertainty. The same six pressure turns that cause Claude Sonnet to flip $71\%$ of uncertain MMLU-Pro questions flip only $10\%$ of its certain ones. This recasts the design of pressure-based sycophancy benchmarks: testing a model on questions where it is already uncertain at baseline overstates its susceptibility to social pressure per se, while testing only on questions where it is confident understates it. Cross-model comparison therefore requires controlling for the baseline uncertainty distribution, which we operationalize as the eligibility criterion of \autoref{sec:methods:setup}; without that control, model-level susceptibility scores conflate \emph{behavioral certainty} with \emph{capitulation propensity}.

\paragraph{Confidence does not track correctness under pressure.} \autoref{tab:rq4} shows that across all five Claude/GPT models on MMLU-Pro and GPQA-Diamond, calibrated correctness drops $25$--$40$ pp from $T_0$ to $T_6$ while mean calibrated confidence stays nearly
flat ($\pm$1--3 pp). The overconfidence gap (Conf.\ $-$ Maj.\ Correct) correspondingly widens from $0.05$--$0.33$ at baseline to $0.21$--$0.88$ at maximum pressure. The most extreme case is Claude Haiku on GPQA-Diamond, where confidence \emph{rises} from $0.805$ to $0.890$ while correctness collapses from $0.679$ to $0.013$. This dissociation generalizes the finding of \citet{mei2025reasoninguncertainty} that
incorrect responses can receive self-verbalized confidence above $85\%$, and extends it to the social-pressure setting: the discriminating signal in a reasoning trace is the divergence between calibrated belief and expressed confidence, not the confidence trajectory in isolation. A confidence-threshold abstention rule built on raw self-reported confidence would catch essentially none of the pressure-induced flips identified by our framework.

\paragraph{Connection to prior multi-turn benchmarks.} TRUTH DECAY \citep{liu2025truthdecayquantifyingmultiturn} and SYCON Bench \cite{hong2025syconbench} establish that multi-turn sycophancy is a distinct failure mode from single-turn agreement, and both report that flip rates change over turns. Our contribution is orthogonal: rather than asking \emph{whether} models flip more under multi-turn pressure, we ask \emph{which questions} are most susceptible and \emph{which models} convert that susceptibility into different sustained-pressure trajectories. The per-turn slope and the pre-pressure uncertainty signal are two black-box quantities that can be added to any existing multi-turn benchmark to yield this finer-grained characterization.
